% Curriculum Vitae -- Yiannis Koutis
% Self-contained; compiles with xelatex / lualatex / tectonic.
\documentclass[11pt,letterpaper]{article}

\usepackage[letterpaper,left=0.85in,right=0.85in,top=0.8in,bottom=0.85in,headsep=16pt,footskip=28pt]{geometry}
\usepackage{amsmath}
\usepackage{libertinus}
\usepackage{xcolor}
\usepackage[explicit]{titlesec}
\usepackage{enumitem}
\usepackage{fancyhdr}
\usepackage{lastpage}
\usepackage[hidelinks]{hyperref}
\usepackage{microtype}

% ---------------------------------------------------------------- colors
\definecolor{accent}{HTML}{1F4E79}
\definecolor{accentlight}{HTML}{E4ECF5}
\definecolor{muted}{HTML}{6B7280}
\definecolor{rulec}{HTML}{B8C4D2}
\definecolor{grayfill}{HTML}{EEF0F3}

% ---------------------------------------------------------------- spacing
\setlength{\parindent}{0pt}
\setlength{\parskip}{0pt}
\widowpenalty=10000
\clubpenalty=10000
\raggedbottom

% ---------------------------------------------------------------- headings
\titleformat{\section}{\large\scshape\color{accent}}{}{0pt}
  {#1\hspace{0.7em}{\color{rulec}\leaders\hrule height 3.6pt depth -3.0pt\hfill}}
\titlespacing*{\section}{0pt}{15pt}{7pt}

\titleformat{\subsection}{\normalsize\bfseries}{}{0pt}{#1}
\titlespacing*{\subsection}{\datecol}{7pt}{3pt}

% ---------------------------------------------------------------- entries
% Every entry hangs from a fixed left column holding a date or label.
\newlength{\datecol}\setlength{\datecol}{1.06in}
\setlength{\emergencystretch}{1em}
\newlength{\entrysep}\setlength{\entrysep}{4pt}

\newcommand{\datefont}{\small\color{muted}}
\newcommand{\entry}[2]{%
  \par\noindent\hangindent=\datecol\hangafter=1
  \makebox[\datecol][l]{\datefont #1}#2\par\vspace{\entrysep}}
% tight variant, for table-like runs (teaching)
\newcommand{\tentry}[2]{{\setlength{\entrysep}{1.5pt}\entry{#1}{#2}}}

% right-aligned trailing text that drops to its own line if needed
\newcommand{\rightpart}[1]{{\unskip\nobreak\hfil\penalty50\hskip1.5em\hbox{}\nobreak\hfil
  #1\parfillskip=0pt\finalhyphendemerits=0\par}}

\newcommand{\badge}[1]{\,{\setlength{\fboxsep}{1.4pt}%
  \colorbox{accentlight}{\textcolor{accent}{\footnotesize\textsc{#1}}}}}
\newcommand{\graybadge}[1]{\,{\setlength{\fboxsep}{1.4pt}%
  \colorbox{grayfill}{\textcolor{muted}{\footnotesize\textsc{#1}}}}}

\newcommand{\note}[1]{\textcolor{muted}{#1}}
\newcommand{\me}{I.~Koutis}
% publication: number, authors, title, venue. The paragraph is left open so
% trailing text (pages, year, badges) joins it; the next \pub or \section closes it.
\newcommand{\pub}[4]{\par\addvspace{\entrysep}\noindent\hangindent=\datecol\hangafter=1
  \makebox[\datecol][l]{\datefont [#1]}#2. #3. \textit{#4}}
% grant: years, title, amount in a fixed right column
\newcommand{\grant}[3]{\entry{#1}{\parbox[t]{\dimexpr\linewidth-\datecol-0.8in}{\raggedright #2}%
  \hfill\makebox[0.65in][r]{#3}}}
\hyphenation{Xiang}

\setlist[itemize]{leftmargin=\dimexpr\datecol+1.1em\relax,labelsep=0.5em,
  itemsep=1.5pt,topsep=2pt,label={\color{muted}\textbullet}}

% ---------------------------------------------------------------- page style
\pagestyle{fancy}
\fancyhf{}
\fancyhead[L]{\small\color{muted}Yiannis Koutis \,·\, Curriculum Vitae}
\fancyfoot[C]{\small\color{muted}\thepage\ of \pageref*{LastPage}}
\renewcommand{\headrulewidth}{0pt}
\fancypagestyle{first}{\fancyhf{}\fancyfoot[C]{\small\color{muted}\thepage\ of \pageref*{LastPage}}}

\begin{document}
\thispagestyle{first}

% ================================================================ header
{\fontsize{26}{30}\selectfont\color{accent}Yiannis Koutis}\par
\vspace{5pt}
{\large Professor of Computer Science}\hfill
\href{mailto:ikoutis@njit.edu}{ikoutis@njit.edu}\par
\vspace{1pt}
New Jersey Institute of Technology\hfill
\href{https://ikoutis.github.io/}{ikoutis.github.io}\par
\vspace{6pt}
{\color{accent}\rule{\linewidth}{0.8pt}}
\vspace{-6pt}

% ================================================================ education
\section{Education}

\entry{2007}{\textbf{PhD}, Carnegie Mellon University, Pittsburgh\\
  Thesis: \textit{Combinatorial and algebraic tools for optimal multilevel algorithms}\\
  Advisor: Gary L. Miller}

\entry{1998}{\textbf{BSc/Diploma}, Computer Engineering and Informatics Department, University of Patras\\
  \textit{Engineering School Valedictorian}\\
  Thesis: \textit{Parallel algorithms for the computation of pseudospectra}\\
  Advisor: Efstratios Gallopoulos}

% ================================================================ appointments
\section{Appointments}

\subsection{Current}
\entry{2026--present}{\textbf{Professor}, New Jersey Institute of Technology}
\entry{2022--present}{\textbf{Associate Chair of Graduate Studies}, Department of Computer Science, NJIT}

\subsection{Past}
\entry{2017--2026}{\textbf{Associate Professor}, New Jersey Institute of Technology}
\entry{2010--2020}{\textbf{Adjunct Faculty}, Carnegie Mellon University}
\entry{2014--2018}{\textbf{Associate Professor}, University of Puerto Rico--Rio Piedras}
\entry{2010--2014}{\textbf{Assistant Professor}, University of Puerto Rico--Rio Piedras}
\entry{2008--2010}{\textbf{Systems Scientist}, Special Faculty, Carnegie Mellon University}
\entry{2007--2008}{\textbf{Postdoctoral Researcher}, Carnegie Mellon University}

\subsection{Visiting}
\entry{Sep--Dec 2014}{\textbf{Visiting Scientist}, Simons Institute}
\entry{Feb--May~2014}{\textbf{Visiting Professor}, ICERM, Brown University}

% ================================================================ awards
\section{Awards and Honors}

\entry{2022}{William J. McCalla ICCAD Best Paper Award (back end)}
\entry{2020}{Ying Wu College of Computing Excellence in Teaching Award}
\entry{2017}{ICALP -- Track A: Best Paper Award}
\entry{2012}{NSF CAREER Award}
\entry{2002}{SIAM Student Travel Award}
\entry{1999}{Best Engineering Student Award, Technical Chamber of Greece}
\entry{1998}{University of Patras Engineering School Valedictorian}
\entry{1993--1998}{Yearly awards from the Greek State Scholarships Foundation}

% ================================================================ grants
\section{Grants}

\grant{2023--2025}{NSF: Cybertraining: Pilot: Cyberinfrastructure-Enabled Machine Learning for
  Understanding and Forecasting Space Weather \badge{co-PI}}{\$190K}
\grant{2020--2022}{NSF: Spectral Network Alignment}{\$150K}
\grant{2018--2021}{NSF: Practice-Friendly Theory and Algorithms for Linear Regression Problems}{\$250K}
\grant{2012--2018}{NSF CAREER: Fast algorithms via a spectral theory for graphs with a prescribed
  cut structure}{\$500K}
\grant{2012--2014}{UPR -- Institute of Functional Nanomaterials (IFN): Computational analysis of
  neural images via spectral methods}{\$50K}
\grant{2010--2013}{NSF: Algorithm Design Using Spectral Graph Theory \badge{co-PI}}{\$500K}
\grant{2008--2009}{University of Pittsburgh Medical Center (UPMC) grant for the development of new
  medical imaging methods based on spectral approaches \badge{co-investigator}}{\$700K}

% ================================================================ conference pubs
\section{Refereed Conference and Workshop Publications}

\pub{31}{M.~Dindoost, O.~Alvarado Rodriguez, B.~Bryg, \me, D.~A.~Bader}
  {HiperMotif: Novel parallel subgraph isomorphism in large-scale property graphs}
  {IEEE High Performance Extreme Computing, HPEC 2025}
\pub{30}{A.~Moradi Karkaj, M.~J.~Nelson, \me, A.~Hoover}
  {Prompt Wrangling: On replication and generalization in large language models for PCG levels}
  {Proceedings of the 15th Workshop on Procedural Content Generation, FDG 2024}\graybadge{workshop}
\pub{29}{I.~Bustany, G.~Gasparyan, A.~Kahng, \me, B.~Pramanik, Z.~Wang}
  {TritonPart: An Open-Source Constraints-Driven General Partitioning Multi-Tool for VLSI Physical Design}
  {42nd International Conference on Computer-Aided Design, ICCAD 2023}
\pub{28}{\me, M.~Wlodarczyk, M.~Zehavi}
  {Sidestepping Barriers for Dominating Set in Parameterized Complexity}
  {18th International Symposium on Parameterized and Exact Computation, IPEC 2023}
\pub{27}{E.~Beikihassan, A.~Parviz, N.~Aghaieabiane, A.~Hoover, \me}
  {Resource-constrained knowledge diffusion processes inspired by human peer learning}
  {26th European Conference on Artificial Intelligence, ECAI 2023}\badge{full oral presentation}
\pub{26}{I.~Bustany, A.~Kahng, \me, B.~Pramanik, Z.~Wang}
  {SpecPart: A supervised spectral framework for hypergraph partitioning solution improvement}
  {41st International Conference on Computer-Aided Design, ICCAD 2022}\badge{best paper award}
\pub{25}{E.~Beikihassan, A.~Hoover, \me, A.~Parviz}
  {Ensemble Learning as a Peer Process}
  {Agent Learning in Open-Endedness Workshop (ALOE), ICLR 2022}\graybadge{workshop}
\pub{24}{\me, B.~Pramanik}
  {Spectral Hypergraph Partitioning Revisited}
  {SIAM Conference on Applied and Computational Discrete Algorithms, ACDA 2021}
\pub{23}{D.~Wei, \me, S.~Basu-Roy}
  {Peer Learning Through Targeted Dynamic Groups Formation}
  {37th IEEE International Conference on Data Engineering, ICDE 2021}
\pub{22}{\me, H.~Le}
  {Spectral Graph Modification for Improved Spectral Clustering}
  {33rd Conference on Neural Information Processing Systems, NeurIPS 2019}
\pub{21}{D.~Calandriello, \me, A.~Lazaric, M.~Valko}
  {Improved Large-Scale Graph Learning through Ridge Spectral Sparsification}
  {35th International Conference on Machine Learning, ICML 2018}
\pub{20}{A.~Kolla, \me, V.~Madan, A.~K.~Sinop}
  {Spectrally Robust Graph Isomorphism}
  {45th International Colloquium on Automata, Languages, and Programming, ICALP 2018}
\pub{19}{A.~Björklund, P.~Kaski, \me}
  {Directed Hamiltonicity and Out-Branchings via Generalized Laplacians}
  {44th International Colloquium on Automata, Languages, and Programming, ICALP 2017}\badge{best paper award}
\pub{18}{I.~Abraham, D.~Durfee, \me, S.~Krinninger, R.~Peng}
  {On fully dynamic graph sparsifiers}
  {57th Annual IEEE Symposium on Foundations of Computer Science, FOCS 2016}
\pub{17}{M.~Cucuringu, \me, S.~Chawla, G.~Miller, R.~Peng}
  {Simple and Scalable Constrained Clustering: A Generalized Spectral Method}
  {19th International Conference on Artificial Intelligence and Statistics, AISTATS 2016}
\pub{16}{C.~Mavroforakis, R.~Garcia-Lebron, \me, E.~Terzi}
  {Spanning edge centrality: large-scale computation and applications}
  {Proceedings of the 24th International World Wide Web Conference, WWW 2015}
\pub{15}{\me}
  {Simple parallel and distributed algorithms for spectral graph sparsification}
  {Proceedings of the 26th Annual Symposium on Parallelism in Algorithms and Architectures, SPAA 2014}
\pub{14}{L.~Brueggeman, M.~Fellows, R.~Fleischer, M.~Lackner, C.~Komusiewicz, \me, A.~Pfandler, F.~Rosamond}
  {Train Marshalling Is Fixed Parameter Tractable}
  {Proceedings of Fun with Algorithms, FUN 2012}
\pub{13}{\me, A.~Levin, R.~Peng}
  {Improved spectral sparsification and numerical algorithms for SDD matrices}
  {Proceedings of the 29th Annual Symposium on Theoretical Aspects of Computer Science, STACS 2012}
\pub{12}{\me, G.~Miller, R.~Peng}
  {A nearly-$m\log n$ solver for SDD linear systems}
  {Proceedings of the 52nd Annual IEEE Symposium on Foundations of Computer Science, FOCS 2011}
\pub{11}{G.~Blelloch, \me, A.~Gupta, G.~Miller, R.~Peng, K.~Tangwongsan}
  {Near linear-work parallel SDD solvers, low-diameter decomposition and low-stretch subgraphs}
  {Proceedings of the 23rd Annual Symposium on Parallelism in Algorithms and Architectures, SPAA 2011}
\pub{10}{\me, G.~Miller, R.~Peng}
  {Approaching optimality for solving symmetric diagonally dominant systems}
  {Proceedings of the 51st Annual IEEE Symposium on Foundations of Computer Science, FOCS 2010}
\pub{9}{G.~Blelloch, \me, G.~Miller, K.~Tangwongsan}
  {Hierarchical Diagonal Blocking with precision reduction applied to combinatorial multigrid}
  {Proceedings of the 23rd ACM/IEEE Conference on High Performance Computing, SC 2010}
\pub{8}{\me, G.~Miller, D.~Tolliver}
  {Combinatorial preconditioners and multilevel solvers for problems in computer vision and image processing}
  {Proceedings of the 5th International Symposium on Visual Computing, ISVC 2009}
\pub{7}{\me, R.~Williams}
  {Limits and applications of group algebras for parameterized problems}
  {Proceedings of the 35th International Colloquium on Automata, Languages and Programming, ICALP 2009}
\pub{6}{C.~Tsourakakis, P.~Drineas, E.~Michelakis, \me, C.~Faloutsos}
  {Spectral counting of triangles in power-law networks via element-wise sparsification}
  {Proceedings of the 2009 International Conference on Advances in Social Networks Analysis and Mining, ASONAM 2009}
\pub{5}{\me}
  {Faster algebraic algorithms for path and packing problems}
  {Proceedings of the 35th International Colloquium on Automata, Languages and Programming, ICALP 2008}
\pub{4}{\me, G.~Miller}
  {Graph partitioning into isolated, high conductance clusters: theory, computation and applications to preconditioning}
  {Proceedings of the 20th Symposium on Parallelism in Algorithms and Architectures, SPAA 2008}
\pub{3}{\me, G.~L.~Miller}
  {A linear work, $O(n^{1/6})$ time, parallel algorithm for solving planar Laplacians}
  {Proceedings of the 18th ACM-SIAM Symposium on Discrete Algorithms, SODA 2007}
\pub{2}{\me}
  {On the hardness of multivariate integration}
  {Proceedings of the 6th International Workshop on Approximation Algorithms for Combinatorial
   Optimization Problems, APPROX 2003}, pp.~122--128
\pub{1}{C.~Bekas, E.~Kokiopoulou, \me, E.~Gallopoulos}
  {Parallel computation of matrix pseudospectra}
  {Proceedings of the 15th ACM International Conference on Supercomputing, ICS 2001}, pp.~260--269

% ================================================================ journal pubs
\section{Journal Publications}

\pub{18}{S.~Rajendra Patil, A.~Parmanand Pandey, \me, M.~Xu}
  {Hierarchical Mamba Meets Hyperbolic Geometry: A New Paradigm for Structured Language Embeddings}
  {Transactions on Machine Learning Research}, 2026
\pub{17}{N.~Aghaieabiane, \me}
  {SGCP: A spectral self-learning method for clustering genes in co-expression networks}
  {BMC Bioinformatics}, 2024
\pub{16}{I.~Bustany, A.~Kahng, \me, B.~Pramanik, Z.~Wang}
  {K-SpecPart: Supervised Embedding Algorithms and Cut Overlay for Improved Hypergraph Partitioning}
  {IEEE Transactions on Computer-Aided Design of Integrated Circuits and Systems}, 2023
\pub{15}{\me, G.~Miller, R.~Peng}
  {A generalized Cheeger inequality}
  {Linear Algebra and Its Applications}, 2023
\pub{14}{N.~Aghaieabiane, \me}
  {A novel calibration step in gene co-expression network construction}
  {Frontiers in Bioinformatics}, 2021
\pub{13}{\me, S.~C.~Xu}
  {Simple parallel and distributed algorithms for spectral graph sparsification}
  {ACM Transactions on Parallel Computing}, 2016
\pub{12}{\me, R.~Williams}
  {Limits and applications of group algebras for parameterized problems}
  {ACM Transactions on Algorithms (TALG)}, 2016
\pub{11}{\me}
  {Multilinear Monomial Detection}
  {Encyclopedia of Algorithms}, 2016\badge{invited}
\pub{10}{\me, R.~Williams}
  {Algebraic Fingerprints for faster algorithms}
  {Communications of the ACM}, January~2016
\pub{9}{\me, A.~Levin, R.~Peng}
  {Faster spectral sparsification and numerical algorithms for SDD matrices}
  {ACM Transactions on Algorithms (TALG)}, 2015
\pub{8}{\me, G.~Miller, R.~Peng}
  {Approaching optimality for solving symmetric diagonally dominant systems}
  {SIAM Journal on Computing, special issue FOCS 2010}, Vol.~43, No.~1, pp.~337--354, 2014\badge{invited}
\pub{7}{G.~Blelloch, \me, A.~Gupta, G.~Miller, R.~Peng, K.~Tangwongsan}
  {Near linear-work parallel SDD solvers, low-diameter decomposition and low-stretch subgraphs}
  {Theory of Computing Systems}, March~2013\badge{invited}
\pub{6}{\me, G.~Miller, R.~Peng}
  {A fast solver for a class of linear systems}
  {Communications of the ACM}, October~2012\badge{invited}
\pub{5}{\me}
  {Constrained multilinear detection for faster functional motif discovery}
  {Information Processing Letters}, Vol.~112, No.~22, pp.~889--892, 2012
\pub{4}{\me, G.~Miller, D.~Tolliver}
  {Combinatorial preconditioners and multilevel solvers for problems in computer vision and image processing}
  {Computer Vision and Image Understanding}, 115(12), pp.~1638--1646, 2011\badge{invited}
\pub{3}{C.~Tsourakakis, P.~Drineas, E.~Michelakis, \me, C.~Faloutsos}
  {Spectral counting of triangles in power-law networks via element-wise sparsification and triangle-based link recommendation}
  {Social Network Analysis and Mining}, Vol.~1, No.~2, pp.~75--81, 2011
\pub{2}{\me}
  {Parameterized complexity and improved inapproximability for computing the largest $j$-simplex in a $V$-polytope}
  {Information Processing Letters}, 1, No.~1, pp.~8--13, 2006
\pub{1}{\me}
  {A faster parameterized algorithm for set packing}
  {Information Processing Letters}, No.~1, pp.~4--7, 2005

% ================================================================ other pubs
\section{Other Publications and Preprints}

\pub{3}{\me}{Dimensionality restrictions on sums over $\mathbb{Z}_p^d$}{CMU-CS-07-103 Technical Report}, 2007
\pub{2}{\me}{Spectrum through pseudospectrum}{arXiv Report No.\ 0701368}, 2001
\pub{1}{\me, E.~Gallopoulos}{Iterations on domains for computing the matrix pseudospectrum}{Manuscript}, 1999

% ================================================================ patents
\section{Patents}

\entry{2014}{\textit{Method and apparatuses for solving weighted planar graphs}. \me, G.~L.~Miller. US~8711146}
\entry{2013}{\textit{Methods for solving graph Laplacians}. \me, G.~L.~Miller. US~8516029}

% ================================================================ invited talks
\section{Invited Talks}

\entry{2019}{\textit{Spectral Graph Modification for Improved Spectral Clustering}\\
  The New York Colloquium on Algorithms and Complexity, CUNY, New York, NY, November~15}
\entry{2019}{\textit{Pragmatic Ridge Spectral Sparsification for Large-Scale Graph Learning}\\
  DIMACS Workshop on Randomized Numerical Linear Algebra, Statistics, and Optimization,
  Rutgers University, New Brunswick, NJ, September~18}
\entry{2018}{\textit{Spectrally Robust Graph Isomorphism}\\
  Simons Institute, Berkeley, CA, September~25}
\entry{2018}{\textit{Directed Hamiltonicity and Out-Branchings via Generalized Laplacians}\\
  Rutgers Business School, Newark, NJ, February~22}
\entry{2017}{\textit{A survey of provably fast linear system solvers}\\
  Alan Turing Institute, London, UK, November~1}
\entry{2017}{\textit{Improved algebraic algorithms for out-branchings problems}\\
  Randomization in Parameterized Complexity, Schloss Dagstuhl, Germany, January}
\entry{2016}{\textit{On fully dynamic graph sparsifiers}\\
  Algebraic and Spectral Graph Theory, Banff International Research Station, Banff, Canada, August~4}
\entry{2016}{\textit{On fully dynamic graph sparsifiers}\\
  Recent Advances on Randomized Numerical Linear Algebra, NII Shonan Meeting, Shonan Village, Japan, July~26}
\entry{2014}{\textit{Spectral algorithms for graph mining and analysis}\\
  Workshop on Algorithms for Modern Massive Data Sets, MMDS 2014, Berkeley, CA, June~17}
\entry{2014}{\textit{Spectral algorithms for graph mining and analysis}\\
  Workshop on Eigenvectors in graph theory and related problems in numerical linear algebra,
  ICERM, Providence, RI, May~7}
\entry{2014}{\textit{Spectral graph sparsification and fast Laplacian solvers}\\
  Boston University, May~5}
\entry{2014}{\textit{Segmenting neurons in 3D EM images}\\
  IBDR PI workshop, Washington, DC, May~2}
\entry{2013}{\textit{Spectral sparsification of graphs: an overview of theory and practical methods}\\
  Workshop on Large Scale Matrix Analysis and Inference, NIPS 2013, Lake Tahoe, NV, December~9}
\entry{2013}{\textit{Algebraization in parameterized algorithms and complexity}\\
  AMS Special Session on the Mathematical Underpinnings of Multivariate Complexity and Algorithm
  Design, and its frontiers, and the field of Incrementalization, San Diego, CA, January~12}
\entry{2012}{\textit{Laplacian Solvers: Theory and Practice}\\
  Workshop on Randomized Numerical Linear Algebra: Theory and Practice, FOCS 2012,
  New Brunswick, NJ, October~20}
\entry{2012}{\textit{SDD Solvers: Bridging the Gap Between Theory and Practice}\\
  Workshop on Algorithms for Modern Massive Data Sets, MMDS 2012, Stanford, CA, July~13}
\entry{2012}{\textit{Spectral graph theory, matrix sums and near-optimal SDD solvers}\\
  SIAM Applied Linear Algebra Conference, Valencia, Spain, June~18}
\entry{2010}{\textit{The power of group algebras in constrained monomial detection problems}\\
  Seminar on the Exact Complexity of NP-hard Problems, Schloss Dagstuhl, Germany, November~10}
\entry{2010}{\textit{Graph Sparsification p.2: An $O(m\log^2 n)$ algorithm for solving SDD systems}\\
  Microsoft Research, Redmond, WA, September~2}
\entry{2010}{\textit{Graph Sparsification p.1: The Combinatorial Multigrid Solver}\\
  Microsoft Research, Redmond, WA, September~2}
\entry{2009}{\textit{How to make a computer see better?}\\
  University of Puerto Rico, Rio Piedras, September~15}
\entry{2009}{\textit{Fast detection of square-free terms in multivariate polynomials: one algo-stone, many algo-birds}\\
  University of Puerto Rico, Rio Piedras, September~14}
\entry{2009}{\textit{M-matrix systems and solvers in Computer Science}\\
  University of Wyoming, March~30}
\entry{2009}{\textit{Spectral Graph Theory meets Practice: The Combinatorial Multigrid Solver}\\
  Yale University, January~23}
\entry{2002}{\textit{Advances in the Theory and Computation of Pseudospectra}\\
  SIAM 50th Anniversary and Annual Meeting, Matrix Spectra and Pseudospectra Minisymposium,
  Philadelphia, PA, July~8--12}

% ================================================================ presentations
\section{Conference and Workshop Presentations}

\entry{2016}{Theoretical Foundations of Data Science: Algorithmic, Mathematical, and Statistical (TFoDS),
  April~28--30}
\entry{2013}{\textit{Fast SDD solvers via sampling by approximate leverage scores}\\
  6th International ERCIM Conference on Computational and Methodological Statistics}
\entry{2011}{\textit{Spectral Algorithms for Segmenting Neurons in their Three-dimensional Space}\\
  Neuroscience 2011 \note{· with R.~Garcia, J.~Farrington, J.~Serrano-Velez, E.~Rosa-Molinar}}
\entry{2009}{\textit{The Combinatorial Multigrid Solver}\\
  The 14th Copper Mountain Conference on Multigrid Methods \note{· with Gary L. Miller}}
\entry{2008}{\textit{Unassisted segmentation of multiple retinal layers via spectral rounding}\\
  Association for Research in Vision and Ophthalmology, ARVO 2008 Annual Meeting
  \note{· with D.~Tolliver, H.~Ishikawa, J.~S.~Schuman, G.~L.~Miller}}
\entry{2007}{\textit{A Linear Work, Parallel Algorithm for Solving Planar Laplacians}\\
  Combinatorial Scientific Computing (CSC07), Costa Mesa, CA \note{· with Gary L. Miller}}
\entry{2006}{\textit{Efficiently Solving Linear Systems using Support Tree Preconditioners}\\
  Parallel Processing for Scientific Computing (PP06), San Francisco, CA, February
  \note{· with Gary L. Miller}}
\entry{2001}{\textit{Hermitian Methods for Computing Eigenvalues}\\
  5th IMACS Conference on Iterative Methods in Scientific Computing, May \note{· with E.~Gallopoulos}}
\entry{1999}{\textit{Iterations on Domains for the Computation of Matrix (Pseudo)-Spectrum}\\
  Conference on the Foundations of Computational Mathematics (FOCM), Oxford, July
  \note{· with E.~Gallopoulos}}

% ================================================================ teaching
\section{Teaching}

\entry{Course design}{Online courses for NJIT: Machine Learning (DS 675), Deep Learning (DS 677)}
\vspace{3pt}

\tentry{NJIT}{Machine Learning (CS/DS 675)\rightpart{\note{S20, F20, F21, Su21, F22, S22, Su22, S23, S24, S26}}}
\tentry{}{Deep Learning (CS/DS 677)\rightpart{\note{F21, S21, F22, S23, Su23, F23, S24, Su24, Su25, Su26}}}
\tentry{}{Data Structures and Algorithms (CS 610)\rightpart{\note{F19, Su20, Su21, S26}}}
\tentry{}{Computational Complexity (CS 611)\rightpart{\note{S18, S19, S20, S21}}}
\tentry{}{Advanced Data Structures and Algorithm Design (CS 435)\rightpart{\note{F17, F18}}}
\tentry{}{Introduction to Data Science (CS 301)\rightpart{\note{F26}}}
\tentry{}{Foundations of Computer Science I (CS 241)\rightpart{\note{F26}}}
\vspace{4pt}
\tentry{UPRRP}{Theory of Computability\rightpart{\note{F15, F16}}}
\tentry{}{Fun with programming interview questions\rightpart{\note{S15}}}
\tentry{}{Theory of Computation\rightpart{\note{F13}}}
\tentry{}{Design and Analysis of Algorithms\rightpart{\note{S12, S17}}}
\tentry{}{High Level Programming Languages\rightpart{\note{S12, S15}}}
\tentry{}{Undergraduate Machine Learning\rightpart{\note{F11, F12, F13}}}
\tentry{}{Numerical Analysis\rightpart{\note{S11, S13}}}
\tentry{}{Linear Algebra for Computer Scientists\rightpart{\note{F10}}}
\vspace{4pt}
\tentry{CMU}{Teaching Assistant, Principles of Programming\rightpart{\note{S02}}}
\tentry{}{Teaching Assistant, Formal Languages, Automata and Computation\rightpart{\note{S01}}}
\vspace{4pt}
\tentry{Patras}{Teaching Assistant, Advanced Scientific Computing\rightpart{\note{F99}}}

% ================================================================ advising
\section{Advising}

\subsection{PhD Students}
\entry{Current}{Soroush Vahidi, Azadeh Naderi, Arash Moradi Karkaj}
\entry{2025}{Ali Parviz\rightpart{\note{First position: Visiting Researcher at Google}}}
\entry{2024}{Ehsan Beikihassan\rightpart{\note{First position: CCC Intelligent Solutions}}}
\entry{2023}{Niloofar Aghaieabiane\rightpart{\note{First position: JP Morgan}}}
\entry{2021}{Huong Le\rightpart{\note{First position: Lecturer at NJIT}}}

\subsection{MS Students}
\entry{2025}{Kadir Altunel}
\entry{2014}{Richard Garcia-Lebron}

\subsection{Former Undergraduate Students}
\entry{}{Richard Garcia-Lebron, Jose Farrington, Karlo Martinez, Leonardo Cardona, Carlos Feliciano,
  Idalyn Mirabal, Alejandro Vientos, Alberto Ruiz, Gustavo Gratacos}

\subsection{PhD Committee Member, NJIT}
\entry{Current}{Mehtab Sidhu, Yupeng Xu, Chunhui Xu, Hongyang Zhang, Mohammad Dindoost, Zhibo Ye,
  Swastik Biswas, Zhenduo Wang, Hamideh Sabaei, Yiran Huang}
\entry{2026}{Shen Fan, Haotian Yin, Siqi Jiang}
\entry{2025}{Oliver Alvarado Rodriguez, Fuhuan Li, Jingyi Gu, Haoran Liu}
\entry{2024}{Wenlu Du, Minjuan Zhang}
\entry{2023}{Mojtaba Zaheri, Mahsa Asadi, Hessam Mohammadi, Sepideh Nikookar, Md Moinul Islam, Xiang Lin}
\entry{2022}{Shibo Yao, Cavidan Yakupoglu, Gerges Firas, Yasser Abduallah}
\entry{2021}{Zhihang Hu, Yunzhe Xue, Dong Wei}
\entry{2020}{Xin Yin}
\entry{2019}{Abdulrhman Fahad Aljouie}

\subsection{PhD Committee Member, External}
\entry{2025}{Jingbang Chen (University of Waterloo)}
\entry{2024}{Yixuan He (University of Oxford), Tijn de Vos (University of Salzburg)}
\entry{2017}{Shen Chen Xu (CMU)}

\subsection{MS Committee Member, NJIT}
\entry{2026}{Ritwik Reddy Kolan}
\entry{2022}{Joseph Patchett}
\entry{2020}{Rahul Basu, Sanyamee Patel, Sarvesh Shukla}

% ================================================================ professional service
\section{Professional Service}

\entry{Panels}{NSF Funding Panels (CCF/AF, IIS/RI) \,·\, DoE Funding Panel \,·\, European Research Council (ERC)}
\entry{PC member}{PKDD-Nectar 2014, 2015; WSDM 2016; WWW 2015--2018; WebConf 2019--2022; KDD 2017;
  KDD-Applied Track 2018; LATIN 2018; CIKM 2018; AAAI 2019--2020; ICCAD 2024--2026;
  DAC 2025--2027; SOSFEM 2027; ACDA 2027}
\entry{Referee}{Conferences: STOC, FOCS, SODA, SPAA, STACS, ESA, MFCS\\
  Journals, including: JACM, ACM TALG, JCSS, Algorithmica, TOPC, IPL, SIMAX}
\entry{Organizer}{Workshop ``Electrical Flows, Graph Laplacians and Algorithms: Spectral Theory and
  Beyond'', ICERM, 2014}

% ================================================================ service NJIT
\section{Service at NJIT}

\subsection{Department Leadership}
\entry{2022--present}{\textbf{Associate Chair of Graduate Studies}, Department of Computer Science}
\vspace{-\entrysep}
\begin{itemize}
  \item Established Barclays--NJIT MS in Computer Science partnership program (2025)
  \item Co-designed three BS program concentrations: AI: Learning and Reasoning, AI: Systems and
        Infrastructure, and Cybersecurity and Privacy (2026)
  \item Developed merit-based faculty salary framework, adopted by CS Department (Spring 2024),
        advancing fairness and transparency in compensation decisions
  \item Redesigned department workflows for course assignments, scheduling, and TA allocation
        (2023--2026); developed supporting software adopted college-wide for all scheduling operations
  \item Updated graduate program catalogs and created Certificate-to-MS pathways
  \item Established Graduate Certificate in Hyperscale Computing
\end{itemize}
\vspace{\entrysep}
\entry{2022--present}{\textbf{MS Committee Chair}, Department of Computer Science}

\subsection{Department Committees \& Administration}
\entry{2024--present}{Committee on Lecturers}
\entry{2023--2024}{Co-administered the Data Structures \& Algorithms PhD Qualifying Exam}
\entry{2022--present}{Organizer of Spring Department Retreats (2023, 2024, 2025)}
\entry{2018--2022}{Administered the Computational Complexity PhD Qualifying Exam}
\entry{2022}{CAREER Mentoring Committee Member}
\entry{2019, 2022--24}{Faculty Search Committee Member}

\subsection{University Service}
\entry{2024--2026}{Center for Educational Innovation and Excellence, Advisory Board}
\entry{2024--2026}{Member of the AI Teaching/Learning Group (charged by the Provost)}
\entry{2024--2025}{Member of the NJIT group participating in AAC\&U's Institute for AI, Pedagogy, and the Curriculum}
\entry{2023--2025}{CS representative, Faculty Senate Committee on Academic Assessment}
\entry{2021--2025}{YWCC representative, Faculty Senate Committee on Research, Scholarship and Creative Academic Activity}
\entry{2019--2024}{CS representative, Committee on Information Technology, Library, and Academic Resources (CITLAR)}

% ================================================================ service UPRRP
\section{Service at UPRRP}

\entry{2016--2017}{Department Coordinator for Student Learning Assessment}
\entry{2012--2014}{Member of the committee for the creation of an MS program in Computer Science}
\entry{2011}{Undergraduate Research Coordinator, Computer Science Department}

\end{document}
